% Generated draft: requires mathematical and source review.
\documentclass[11pt]{article}
\usepackage[margin=25mm]{geometry}
\usepackage{fontspec}
\setmainfont{Latin Modern Roman}
\usepackage{amsmath,amssymb,mathtools,unicode-math}
\setmathfont{Latin Modern Math}
\usepackage{enumitem,needspace,xurl,hyperref}
\setlength{\parindent}{0pt}
\setlength{\parskip}{5pt}
\setlength{\emergencystretch}{3em}
\newcommand{\N}{\mathbb{N}}
\newcommand{\Z}{\mathbb{Z}}
\newcommand{\Q}{\mathbb{Q}}
\newcommand{\R}{\mathbb{R}}
\newcommand{\C}{\mathbb{C}}
\newcommand{\F}{\mathbb{F}}
\newcommand{\A}{\mathbb{A}}
\newcommand{\PP}{\mathbb{P}}
\newcommand{\E}{\mathbb{E}}
\newcommand{\eps}{\varepsilon}
\newcommand{\dd}{\,\mathrm{d}}
\begin{document}
\section*{UnsolvedMath ID: 701387}
\textbf{Source problem code:} catalogue entry
\subsection{Definitions and mathematical statement}
A variety of groups is a class \(\mathcal V\) consisting of all groups satisfying a specified collection of group identities. For each set \(X\), let \(F_{\mathcal V}(X)\) be the relatively free group in \(\mathcal V\) on \(X\): every map from \(X\) to a group \(G\in\mathcal V\) extends uniquely to a homomorphism \(F_{\mathcal V}(X)\to G\). Write \(H_n(G;\Z)\) for group homology with trivial integral coefficients. An abelian group \(A\) is torsion-free if \(ma=0\), for \(a\in A\) and an integer \(m\geq 1\), implies \(a=0\).

Suppose that
\[
\text{for every set }X\text{ and every integer }n\geq 0,\qquad H_n(F_{\mathcal V}(X);\Z)\text{ is torsion-free}.
\]
Must every group in \(\mathcal V\) be abelian? Equivalently, must the identity
\[
xy=yx
\]
hold in every group in \(\mathcal V\)? Here the conclusion permits any subvariety of the variety of abelian groups; it does not require equality with the entire variety of abelian groups. The source places no rank restriction on its phrase ``any free group in \(\mathcal V\).''

The supplied background states that free metabelian groups may have \(2\)-torsion in their second integral homology groups, and free nilpotent groups of class \(2\) may have \(3\)-torsion in their fourth integral homology groups. A group is metabelian when its commutator subgroup is abelian. A group has nilpotency class \(2\) when its commutator subgroup is nontrivial and central. The background assertions are retained as supplied qualifications: they do not assert the presence of such torsion for every rank, and the record specifies neither ranks nor separate supporting references.
\subsection{Short English statement}
R. Mikhailov asks whether a variety of groups must consist entirely of abelian groups if every relatively free group in the variety has torsion-free integral homology in every dimension. The supplied record labels this problem open; this label and the accompanying torsion examples have not been independently verified.
\subsection{Sources}
\begin{itemize}
\item Local UnsolvedMath registry, permanent ID \(701387\), exact source problem code \texttt{catalogue entry}, title \emph{Kourovka Notebook Problem 17.85}; supplied upstream revision \texttt{d371ee917ac9647fbb06cd6fc529deec6f3dda82}.
\item \emph{The Kourovka Notebook (Unsolved Problems in Group Theory)}, Problem 17.85, proposed by R. Mikhailov in 2010. The supplied background identifies the 17th issue (2010) and the 21st issue edition (2026), with identifier \texttt{arXiv:1401.0300v48}. The supplied collection description names E. I. Khukhro and V. D. Mazurov as editors.
\item The homological torsion examples and the status label \texttt{open} are supplied by the local record. No separate citations for those examples, literature-verification date, or external source links are supplied. The cited edition and its contents have not been independently checked for this draft.
\end{itemize}
\end{document}
